\chapter*{Abstract}
\addcontentsline{toc}{chapter}{Abstract}

This thesis develops a real-time controller for a seven-degree-of-freedom
collaborative robot intended for surface grinding with a rectangular tool on a
planar surface. Cartesian impedance control is used to generate the required
pressing motion while allowing passive, contact-induced rotation of the tool
towards surface alignment. This behaviour is intended to reduce the effect of
angular mismatch at contact entry, which can otherwise produce edge-first
contact and increased contact loads.

Directional translational and rotational compliance was defined relative to the
surface. In addition, a virtual centre of compliance could be displaced from the
tool centre point to couple translational pressing with rotational motion and
thereby influence the preferred direction of alignment.

Surface-contact experiments investigated the angular-error component
about the first surface tangent relative to the calibrated surface.
Different stiffness and centre-of-compliance settings were compared.
Increasing rotational stiffness increased the angular error by about
\(276\,\%\), whereas the tested tangential translational stiffness variation
changed it by about \(4\,\%\) of the baseline error.
The effect of centre-of-compliance displacement depended on the entry
direction. The lowest measured mean error magnitudes were about
\(33\,\%\) and \(34\,\%\) below their respective tool-centre-point references,
at different tested positions.
The evaluated angular-error component could pass through zero.
Further contact-induced rotation could increase its final magnitude.

A separate pose-hold experiment evaluated the null-space controller. Projected
damping reduced cumulative joint motion by about \(25\,\%\), while
singular-value conditioning reduced the remaining net joint motion by
approximately \(99.8\,\%\) to \(99.9\,\%\).
